%----------------------------------------------------------------------------------------
%	SECTION 5: TRAJECTORY TRACKING VIA MPC
%----------------------------------------------------------------------------------------



% (section 5) task 4 - Trajectory tracking with MPC:
%     - prediction horizon
%     - perturbation
%     - constraint u (max u)    
%
\section{Task 4: Trajectory Tracking via MPC}

While LQR is effective for unconstrained tracking, it \textbf{cannot explicitly handle actuator saturation}. We implemented a linear \textbf{Model Predictive Control} (MPC) scheme to address this.

\subsection{MPC Formulation}
At each time step $t$, we solve a finite-horizon optimization problem:
\begin{align*}
\min_{\Delta u} \quad & \sum_{t=0}^{T_{pred}-1} \left( \|\Delta x_t\|_{Q}^2 + \|\Delta u_t\|_{R}^2 \right) + \|\Delta x_{T_{pred}}\|_{Q_T}^2 \\
\text{s.t.} \quad & \Delta x_{k+1} = A_{t+k} \Delta x_k + B_{t+k} \Delta u_k \\
& u_{min} \leq u^*_{t+k} + \Delta u_k \leq u_{max}
\end{align*}
where $T_{pred}$ is the prediction horizon.

\subsection{Key Features}
\begin{itemize}
\item \textbf{Constraints:} We imposed hard constraints on the input torque ($u \in [u_{min}, u_{max}]$) to \textbf{respect physical motor limits}, which LQR ignores. \\
In particular, we set $u_{max} = \max(|u^*|)$ and $u_{min} = -u_{max}$, so as \textbf{not to exceed the torque needed to complete the nominal trajectory} with no disturbances; managing \\
to find a way through thanks to the prediction window and the optimization.
\item \textbf{Receding Horizon:} Only the first control input is applied, and the optimization is repeated at the next step. Furthermore we opted for a prediction horizon $T_{pred}$\\
that could serve as a \textbf{good compromise between performance and computational load}. In particular, we set $T_{pred} = \min(100, \frac{\text{len}(time)}{2})$.
\item \textbf{Perturbations:} Similar to the LQR case, we introduced random initial state perturbations to evaluate robustness. We opted for bigger values, about $1.5$ times those used in the LQR case.
\item \textbf{Solver:} We utilized CasADi with IPOPT to solve the constrained QP/NLP at each step.
\end{itemize}

\subsection{Simulation Results}
The MPC controller showed superior performance in scenarios where large perturbations would cause an LQR controller to request infeasible torques.\\
Moreover, the ability to handle constraints makes MPC a versatile choice.\\
Down below some results are shown (see Figure \ref{fig:placeholder1}, \ref{fig:placeholder2}, \ref{fig:placeholder3}).

\begin{figure}[]
    \centering
    \includegraphics[width=0.1\linewidth]{figs/placeholder.jpg}
    \caption{Your descriptive caption goes here.}
    \label{fig:placeholder1}
\end{figure}

\begin{figure}[]
    \centering
    \includegraphics[width=0.1\linewidth]{figs/placeholder.jpg}
    \caption{Your descriptive caption goes here.}
    \label{fig:placeholder2}
\end{figure}

\begin{figure}[]
    \centering
    \includegraphics[width=0.1\linewidth]{figs/placeholder.jpg}
    \caption{Your descriptive caption goes here.}
    \label{fig:placeholder3}
\end{figure}

\newpage